
\medskip
\phantomsection\label{prop:sga3-VIB-6.2}%
\noindent\textbf{Proposition 6.2.---}%
\footnote{\label{note:sga3-VIB-51}Editor's note: we have modified
the statement as indicated in editor's \hyperref[note:sga3-VIB-49]{note 49}.}
\emph{Let \(G\) be an \(S\)-group. For a subfunctor \(T\) of the
functor \(G\), consider the following property:}
\[
  (+_{\!f})\qquad
  \text{for every \(s\in S\), \(T_s\) is representable by a
  closed subscheme of \(G_s\).}
\]
\emph{Let \(u,w:X\to G\) and \(v:Y\to G\) be morphisms of
\(S\)-schemes. Then:}
\begin{enumerate}
\item[\textup{(i)}]
  \(\underline{\operatorname{Transp}}(u,w)\) and
  \(\underline{\operatorname{Centr}}(u)
    =\underline{\operatorname{Transp}}(u,u)\)
  satisfy condition \((+_{\!f})\).
\item[\textup{(ii)}]
  \(\underline{\operatorname{Transp}}_G(u(X),v(Y))\) and
  \(\underline{\operatorname{Norm}}_G v(Y)\) satisfy condition
  \((+_{\!f})\) if, for every \(s\in S\), \(v_s\) is a closed
  immersion.
\item[\textup{\SGAInlineRefTarget{prop:sga3-VIB-6.2:item:iii}(iii)}]
  \(\underline{\operatorname{Transp}}^{\mathrm{str}}_G(X,Y)\)
  satisfies condition \((+_{\!f})\) if, for every \(s\in S\),
  \(u_s\) and \(v_s\) are closed immersions.
\end{enumerate}

This follows from
\hyperref[rem:sga3-VIB-6.1.1]{Remark~6.1.1} and
\hyperref[cor:sga3-VIB-6.2.5]{Corollary~6.2.5} below.%
\footnote{\label{note:sga3-VIB-52}Editor's note: the original
referred to the results of Expos\'e VIII, \S6. For the reader's
convenience, we have reproduced these results (with the exception
of VIII, \SGARefLink{sga3:viii:proposition:6:3}{6.3} and \SGARefLink{sga3:viii:remark:6:8}{6.8}) in \SGARefLink{def:sga3-VIB-6.2.1}{6.2.1}--\SGARefLink{cor:sga3-VIB-6.2.5}{6.2.5} below. Moreover, this had
been suggested by A.\ Grothendieck in a note at the beginning of
VIII.6: ``The present section is independent of the theory of
diagonalizable groups; its natural place would be in
VI\(_B\).''}

\medskip
\phantomsection\label{def:sga3-VIB-6.2.1}%
\noindent\textbf{Definition 6.2.1.---}
Let \(f:X\to S\) be a morphism of schemes. One says that \(f\) is
\emph{essentially free}, or that \(X\) is \emph{essentially free
over \(S\)}, if one can find a covering of \(S\) by affine open
subsets \(S_i\), for every \(i\) an affine \(S_i\)-scheme \(S'_i\)
faithfully flat over \(S_i\), and a covering
\((X'_{ij})_j\) of
\[
  X'_i=X\times_S S'_i
\]
by affine open subsets \(X'_{ij}\), such that, for every
\((i,j)\), the ring of \(X'_{ij}\) is a projective%
\footnote{\label{note:sga3-VIB-53}Editor's note: on the one hand,
we have replaced the word ``free'' by ``projective'', as
indicated in VIII, \SGARefLink{sga3:viii:remark:6:8}{Remark 6.8}. On the other hand, the notion of an
essentially free \(S\)-scheme is closely related to the geometric
notion of an \(S\)-scheme that is flat and pure, introduced and
studied by M.\ Raynaud and L.\ Gruson ([RG71]); see the addition
\hyperref[par:sga3-VIB-6.9]{6.9} below.}
module over the ring of \(S'_i\).

\medskip
\phantomsection\label{prop:sga3-VIB-6.2.2}%
\noindent\textbf{Proposition 6.2.2.---}
\begin{enumerate}
\item[\textup{a)}]
  \emph{If \(X\) is essentially free over \(S\), it is flat over
  \(S\); the converse holds if \(S\) is Artinian.}
\item[\textup{b)}]
  \emph{If \(S\) is the spectrum of a field, every \(S\)-scheme
  is essentially free over \(S\).}
\item[\textup{c)}]
  \emph{If \(X\) is essentially free over \(S\), then
  \(X'=X\times_S S'\) is essentially free over \(S'\), for every
  \(S'\to S\). The converse holds if \(S'\to S\) is faithfully
  flat and quasi-compact.}
\end{enumerate}

The proof is immediate, using for the converse in \textup{a)} the
fact that a flat module over an Artinian local ring is free.%
\footnote{\label{note:sga3-VIB-54}Editor's note: indeed, let
\((A,\mathfrak m)\) be an Artinian local ring, \(k\) its residue
field, \(M\) an arbitrary \(A\)-module, and let
\((x_i)_{i\in I}\) be elements of \(M\) whose images form a basis
of \(M/\mathfrak mM\) over \(k\). Let \(F\) be the free
\(A\)-module with basis \((e_i)_{i\in I}\), and let
\(\varphi:F\to M\) be the \(A\)-morphism defined by
\(\varphi(e_i)=x_i\). Then \(Q=\operatorname{Coker}\varphi\)
satisfies \(Q=\mathfrak mQ\), whence, since \(\mathfrak m\) is
nilpotent, \(Q=0\). Suppose moreover that \(M\) is flat over
\(A\); then \(K=\operatorname{Ker}\varphi\) satisfies
\(K\otimes_Ak=0\), that is, \(K=\mathfrak mK\), whence \(K=0\).}

The introduction of
\hyperref[def:sga3-VIB-6.2.1]{Definition~6.2.1} is justified by
the following theorem.
